
\section{三种滑模设计方法对比}

三种方法采用同一周期滑模控制律
\[
    u_{i\ell}=R_{i\ell}\sgn[\sin(\pi s_{i\ell}/\varepsilon)],\qquad
    R_{i\ell}=\frac{|\Delta_{i\ell}|+\bar w_{i\ell}+\varepsilon/T^{\mathrm r}_{i\ell}}{\ub_{i\ell}}.
\]
其中 $\bar w_{i\ell}$ 是逐分量扰动界，$T^{\mathrm r}_{i\ell}$ 是期望的单个周期单元到达时间；区别仅在于滑模变量（即 $\Delta_i$）的定义。详细推导见 result1/2/3。

\begin{table}[H]
\centering
\caption{result1 / result2 / result3 对比}
\begin{tabular}{>{\raggedright\arraybackslash}p{2.2cm}|>{\raggedright\arraybackslash}p{3.9cm}|>{\raggedright\arraybackslash}p{3.9cm}|>{\raggedright\arraybackslash}p{3.9cm}}
\hline
& 方法1 & 方法2（积分滑模） & 方法3（滤波排斥项）\\
\hline
滑模变量 &
$s_{i0}=v_{i0}+k_1p_{i0}-\varphi_i+k_2\eta_i$ &
$s_{i0}=\eta_i+k_2p_{i0}+v_{i0}$ &
$s_{i0}=v_{i0}+k_1p_{i0}-\Phi_i+k_2\eta_i$\\
\hline
$\dot\eta_i$ &
$k_1p_{i0}-\varphi_i$ &
$k_1p_{i0}-\varphi_i$ &
$k_1p_{i0}-\Phi_i$\\
\hline
$\Delta_i$ &
$-\dot\varphi_i+k_1v_{i0}+k_2(k_1p_{i0}-\varphi_i)$ &
$k_1p_{i0}-\varphi_i+k_2v_{i0}$ &
$-\dot\Phi_i+k_1v_{i0}+k_2(k_1p_{i0}-\Phi_i)$\\
\hline
需邻居速度 &
是（$\dot\varphi_i$） &
否 &
否\\
\hline
降阶动力学 &
含 $\dot\varphi_i$，非标准 &
$\ddot p_{i0}+k_2\dot p_{i0}+k_1p_{i0}=\varphi_i$ &
同方法1（用 $\Phi_i$）\\
\hline
P1/P2 证明 &
自证（障壁势函数） &
自证（P1）；P2 需额外 remark（\(\dot s_i\) 有界） &
自证（增广势函数；P2 需快滤波备注）\\
\hline
P3 速度一致 &
未证 &
自证（Barbalat） &
未证\\
\hline
\end{tabular}
\end{table}

要点：方法1 是基线，但滑模变量直接含 $-\varphi_i$，求导后出现 $\dot\varphi_i$，需要邻居相对速度；方法2 用积分滑模把 $\varphi_i$ 关进积分器，既无需邻居速度，又得到标准名义合围动力学；方法3 用一阶滤波 $\dot\Phi_i=-\gamma_2\Phi_i+\varphi_i$ 替代 $\varphi_i$，同样去掉邻居速度依赖且保持方法1 形式，但 P2 需要额外快滤波条件。三者到达性证明与关键不等式 $|b|R>|N|$ 结构完全一致。
